% Fiche C++ — recto-verso A4 paysage
% Source : cours POO 3A SEI (L01–L03, POO-1/2/3.md, labs) — contenu repris, corrigé et complété.
% Complément de la fiche C (../C) : types, opérateurs, contrôle, pointeurs, mémoire,
% préprocesseur et compilation n'y sont PAS répétés ; seul ce qu'ajoute le C++ figure ici.
% Compiler : make (ici ou à la racine), ou xelatex cpp_cheat.tex — une seule passe suffit
% Apparence : ../librairie (fiche.sty + theme-cpp.tex). Macros de contenu :
%   colonne[titre][couleur] (3 par page, séparées par \colbreak) · bloc{titre}{cat}
%   catégories : syn typ ctl ex io lib (couleurs et icônes dans le thème)
%   code · \V|...| code en ligne · «...» dans le code = à remplir
%   \sub sous-titre · \note remarque · \piege piège · \badge{TXT} · schema
%   optab · tab3 · compare · grille{n} (largeurs mesurées par la librairie)
% Le code est en police à chasse fixe : les commentaires sont alignés à la main
% (attention : les « » ne sont pas affichés, ils ne comptent pas dans l'alignement).
% Après compilation : grep "\[fiche\]" cpp_cheat.log (place libre, collisions).
\documentclass{article}
\usepackage[cpp]{../librairie/fiche}   % apparence : ../librairie (fiche.sty + theme-cpp.tex)
% --- Langage C++ pour listings ---
% Classes de mots-clés : [1] rose · [2] cyan · [3] vert · [4] couleur du thème
% (styles dans ../librairie/fiche.sty). Chargé par cpp_cheat.tex.
\lstdefinelanguage{CODE}{
  sensitive=true, alsoletter={\#},
  morekeywords=[1]{if,else,for,while,do,return,break,continue,switch,case,default,
    class,struct,public,private,protected,virtual,override,final,template,typename,
    namespace,using,new,delete,this,friend,operator,try,catch,throw,noexcept,explicit,
    static,const,constexpr,mutable,inline,enum,typedef,nullptr,true,false,sizeof,
    auto,extern,static_cast,dynamic_cast,const_cast,reinterpret_cast},
  morekeywords=[2]{int,float,double,char,void,bool,short,long,unsigned,signed,size_t,
    uint8_t,uint32_t,string,vector,list,map,set,unordered_map,array,deque,pair,
    function,unique_ptr,shared_ptr,ostream,istream,ifstream,ofstream,istringstream,
    exception,runtime_error,invalid_argument,out_of_range,bad_alloc,iterator},
  morekeywords=[3]{std,cout,cin,cerr,endl,getline,make_unique,make_shared,move,
    sort,find,find_if,count_if,accumulate,remove_if,for_each,copy,to_string,
    stoi,stod,setw,setprecision,fixed,
    \#include,\#pragma,\#define,\#ifndef,\#endif},
  morecomment=[l]{//}, morecomment=[s]{/*}{*/}, morestring=[b]", morestring=[b]',
}
   % coloration syntaxique : \lstdefinelanguage{CODE}

\begin{document}
%%%%%%%%%%%%%%%%%%%%%%%%%%%%%%%%%%%%%%%%%%%%%%%%%%%%%%%%%%%%%%%%%%%%%%%%%%%%%%
%                     RECTO : LE LANGAGE OBJET
%%%%%%%%%%%%%%%%%%%%%%%%%%%%%%%%%%%%%%%%%%%%%%%%%%%%%%%%%%%%%%%%%%%%%%%%%%%%%%
\entete{C++}{recto · ce que le C++ ajoute au C}
\noindent
% ======================= colonne 1 : du C au C++ =======================
\begin{colonne}[du C au C++][col-a]
\begin{bloc}{Du C au C++}{syn}
\begin{code}
#include <iostream>           // sans .h
bool ok = true;               // type natif
int *p = nullptr;             // remplace NULL
auto x = 3.0;                 // type déduit : double
for (int v : tab) { ... }     // parcours de conteneur
for (auto &v : tab) v *= 2;   // par réf. : modifie
int f(int a, int b = 0);      // argument par défaut
\end{code}
\sub{Surcharge de fonctions}
\begin{code}
void print(int v);            // même nom,
void print(double v);         // paramètres différents
void print(const char *s);
int  print(int v);            // ERREUR : retour seul
\end{code}
\note{Signature = nom + portée + paramètres (+ \V|const|), \textbf{pas le retour} ;
encodée dans le symbole : \V|print(int)| $\to$ \V|_Z5printi|. \V|extern "C"| pour lier du C.}
\end{bloc}

\begin{bloc}{Références}{typ}
\begin{code}
int i = 42;
int &r = i;                   // alias, à initialiser
r = 10;                       // i vaut 10 (pas de *)
void inc(int &x) { x++; }     // paramètre de sortie
void show(const Big &b);      // entrée sans copie
inc(i);                       // appel sans &
\end{code}
\begin{compare}
           & \V|T *p| pointeur  & \V|T &r| référence \\ \hline
Init.      & facultative        & \textbf{obligatoire} \\
Nul        & \V|nullptr|        & impossible \\
Réaff.     & \V|p = &autre;|    & non : \V|r = b| copie b \\
Arithm.    & \V|p++|  \V|p + i| & non \\
Accès      & \V|*p|  \V|p->m|   & \V|r|  \V|r.m|
\end{compare}
\piege{\V|f(int)| et \V|f(int &)| ensemble : \V|f(i)| est ambigu.
Ne jamais renvoyer une référence vers une variable locale.}
\end{bloc}

\begin{bloc}{Les 4 casts}{typ}
\begin{tab3}[0mm]{33mm}
 & \V|static_cast<T>(x)|      & vérifié à la compilation \\
 & \V|dynamic_cast<D *>(b)|   & à l'exécution ; sinon \V|nullptr| \\
 & \V|const_cast<T *>(p)|     & ajoute / retire \V|const| \\
 & \V|reinterpret_cast<T>(x)| & relit les bits bruts \\
 & \V|(T)x|                   & cast C : à éviter
\end{tab3}
\begin{code}
Base *b = new Derived();
if (Derived *d = dynamic_cast<Derived *>(b))
    d->execute();             // b est bien un Derived
auto *r = reinterpret_cast<Regs *>(0x40001000);
\end{code}
\note{\V|dynamic_cast| : classe polymorphe (une méthode \V|virtual|) exigée.}
\end{bloc}

\begin{bloc}{enum class et namespace}{syn}
\begin{code}
enum class Color : uint8_t { Red, Green, Blue };
Color c = Color::Red;         // portée obligatoire
int i = static_cast<int>(c);  // conversion explicite
namespace phelma::proj1 { void print(); }  // C++17
phelma::proj1::print();
namespace pp = phelma::proj1; // alias
using namespace std;          // jamais dans un .hpp
\end{code}
\end{bloc}
\end{colonne}\colbreak
% ======================= colonne 2 : classes =======================
\begin{colonne}[classes][col-b]
\begin{bloc}{Anatomie d'une classe}{ctl}
\begin{code}
// Circle.hpp : interface
#pragma once                  // ou garde #ifndef
class Circle {
public:
    Circle(double r = 1.0);   // constructeur
    ~Circle();                // destructeur
    double area() const;      // ne modifie pas *this
    static int count();       // sans objet, sans this
private:
    double radius;            // attribut
    static int nb;            // un seul pour tous
};
// Circle.cpp : implémentation
#include "Circle.hpp"
int Circle::nb = 0;           // définition du static
Circle::Circle(double r) : radius(r) { nb++; }
Circle::~Circle() { nb--; }
double Circle::area() const {
    return 3.14159 * radius * radius;
}
// main.cpp
Circle c(2.0);                // pile : auto-détruit
Circle *pc = new Circle(3);   // tas
pc->area();
delete pc;                    // destr. + libération
Circle::count();
\end{code}
\note{Accès vérifié à la compilation seulement. Un objet ne contient que ses attributs ;
les méthodes reçoivent \V|this| en paramètre caché.}
\end{bloc}

\begin{bloc}{Construction et destruction}{ctl}
\begin{code}
class Duration {
public:
    Duration();                   // par défaut
    explicit Duration(long s);    // pas de conversion
    Duration(const Duration &o);  // copie
    Duration &operator=(const Duration &o);
    ~Duration();
private:
    long value;                   // pas const (operator=)
    std::string unit;
};
Duration::Duration(long s)
    : value(s), unit("s") {}      // liste d'init.
Duration d1(60);                  // constructeur
Duration d2 = d1;                 // constr. de copie
d2 = d1;                          // operator=
void f(Duration d);               // copiée à l'appel
Duration d3 = 60;                 // ERREUR : explicit
\end{code}
\begin{optab}[15mm]
Liste init.  & obligatoire pour : membres \V|const|, références, membres sans constructeur
               par défaut, classe de base à paramètres \\
Ordre        & celui de \textbf{déclaration} des membres, pas celui de la liste \\
Cycle        & base $\to$ membres $\to$ corps ; destruction en ordre inverse \\
Implicites   & copie, \V|operator=|, destructeur ; \V|T()| seulement si aucun constructeur.
               + déplacement (C++11). \V|= default| les demande, \V|= delete| les interdit
\end{optab}
\end{bloc}
\end{colonne}\colbreak
% ======================= colonne 3 : héritage =======================
\begin{colonne}[héritage][col-c]
\begin{bloc}{Héritage et polymorphisme}{ctl}
\begin{code}
class Shape {                        // abstraite
public:
    Shape(std::string c) : color(c) {}
    virtual ~Shape() = default;      // OBLIGATOIRE
    virtual void draw() const = 0;   // virtuelle pure
    virtual double area() const { return 0; }
protected:
    std::string color;               // pour les filles
};
class Rect : public Shape {
public:
    Rect(int w, int h) : Shape("red"), w(w), h(h) {}
    void draw() const override;      // redéfinition
    double area() const override { return w * h; }
private:
    int w, h;
};
void render(const Shape &s) {        // upcast
    s.draw();                        // -> Rect::draw()
}
std::vector<Shape *> v = {new Rect(1, 2)};
for (Shape *s : v) s->draw();        // appel dynamique
\end{code}
\begin{optab}[14mm]
\V|virtual|  & résolu à l'exécution, selon le type réel \\
\V|= 0|      & virtuelle pure $\Rightarrow$ classe abstraite, non instanciable ;
               que des virtuelles pures = interface \\
\V|override| & le compilateur vérifie la redéfinition · \V|final| l'interdit
\end{optab}
\piege{sans destructeur \V|virtual|, \V|delete| d'un \V|Shape *| oublie la partie \V|Rect|.
passer une base \textbf{par valeur} (\V|void f(Base b)|) tronque l'objet dérivé (\emph{slicing}) : passer \V|const Base &|.}
\end{bloc}

\begin{bloc}{vtable et coût}{typ}
\begin{code}
struct A { uint32_t id, flags; void f(); };
struct B { uint32_t id, flags; virtual void f(); };
sizeof(A);                    // 8
sizeof(B);                    // 16 : + vptr (8 octets)
\end{code}
\note{Une \textbf{vtable} (adresses des méthodes virtuelles) par classe, un \textbf{vptr}
par objet. Appel virtuel = une indirection.}
\end{bloc}

\begin{bloc}{Accès et héritage}{ctl}
\begin{tab3}[9mm]{19mm}
Accès  & \V|public|    & tout le monde \\
       & \V|protected| & la classe et ses dérivées \\
       & \V|private|   & la classe (et ses \V|friend|) \\
Défaut & \V|struct|    & tout \V|public| \\
       & \V|class|     & tout \V|private| ; héritage \V|private| par défaut
\end{tab3}

\begin{optab}[17mm]
Composition & « possède un » : attribut (\V|Movie| a une \V|Duration|) \\
Héritage    & « est un » : \V|Rect| est une \V|Shape| (substituable)
\end{optab}
\begin{tab3}[9mm]{21mm}
Mode & \V|: public B|    & droits inchangés \\
     & \V|: protected B| & public devient protected \\
     & \V|: private B|   & tout devient private (défaut)
\end{tab3}
\begin{code}
class Product { public: float price; };
class Paper : public virtual Product { ... };
class Elec  : public virtual Product { ... };
class Smartpad : public Paper, public Elec {};
pad.price = 9.5f;             // une seule copie de price
\end{code}
\note{Sans \V|virtual| : deux \V|Product| (diamant), il faut écrire \V|pad.Paper::price|.
Usage typique : hériter de plusieurs interfaces.}
\end{bloc}
\end{colonne}

%%%%%%%%%%%%%%%%%%%%%%%%%%%%%%%%%%%%%%%%%%%%%%%%%%%%%%%%%%%%%%%%%%%%%%%%%%%%%%
%                     VERSO : OUTILS ET BIBLIOTHÈQUE
%%%%%%%%%%%%%%%%%%%%%%%%%%%%%%%%%%%%%%%%%%%%%%%%%%%%%%%%%%%%%%%%%%%%%%%%%%%%%%
\newpage
\entete{C++}{verso · ressources, généricité et bibliothèque}
\noindent
% ======================= colonne 1 : opérateurs, ressources =======================
\begin{colonne}[opérateurs · RAII][col-a]
\begin{bloc}{Surcharge d'opérateurs}{syn}
\begin{code}
class Complex {
public:
    Complex(double re = 0, double im = 0)
        : re(re), im(im) {}
    Complex operator+(const Complex &o) const {
        return Complex(re + o.re, im + o.im);
    }
    Complex &operator+=(const Complex &o) {
        re += o.re;  im += o.im;
        return *this;                // chaînage
    }
    Complex &operator++();           // ++c
    Complex operator++(int);         // c++ (int fictif)
    friend std::ostream &operator<<(std::ostream &os,
                                    const Complex &c);
private:
    double re, im;
};
std::ostream &operator<<(std::ostream &os,
                         const Complex &c) {
    return os << c.re << " + " << c.im << "i";
}
Complex a(1, 2), b = a + 3.0;        // 3.0 -> Complex
std::cout << a << '\n';
\end{code}
\note{\textbf{Méthode} si l'opérande gauche est l'objet, sinon \textbf{fonction libre}
(\V|cout << c|, \V|2.0 + c|). Interdits : \V|::  .  .*  ?:  sizeof|.}
\end{bloc}

\begin{bloc}{Allocation et ressources}{lib}
\begin{code}
T *p = new T(args);                  // + constructeur
delete p;                            // + destructeur
T *t = new T[n];                     // n fois T()
delete[] t;                          // pas delete seul
auto u = std::make_unique<T>(args);  // delete auto
auto s = std::make_shared<T>(args);  // partagé
\end{code}
\note{Échec : exception \V|std::bad_alloc|. Pointeurs intelligents : \V|<memory>|.}
\sub{Règle des trois}
\begin{code}
class Buf {                          // possède un new[]
    int *data;
    size_t n;
public:
    Buf(size_t n) : data(new int[n]), n(n) {}
    ~Buf() { delete[] data; }        // 1. destructeur
    Buf(const Buf &o)                // 2. copie profonde
        : data(new int[o.n]), n(o.n) {
        std::copy(o.data, o.data + n, data);
    }
    Buf &operator=(const Buf &o) {   // 3. affectation
        if (this != &o) {            // auto-affectation
            delete[] data;
            n = o.n;
            data = new int[n];
            std::copy(o.data, o.data + n, data);
        }
        return *this;
    }
};
\end{code}
\note{Copie par défaut = superficielle : deux objets, un seul tableau $\Rightarrow$ double
\V|delete[]|. \textbf{RAII} : ressource prise dans le constructeur, rendue dans le destructeur.}
\end{bloc}
\end{colonne}\colbreak
% ======================= colonne 2 : généricité =======================
\begin{colonne}[généricité · lambdas][col-b]
\begin{bloc}{Templates}{typ}
\begin{code}
template <typename T>
T maxOf(T a, T b) { return a > b ? a : b; }
maxOf(3, 7);                  // T = int, déduit
maxOf<double>(3, 7.5);        // T imposé

template <typename T1, typename T2>
class Pair {
public:
    Pair(T1 a, T2 b) : first(a), second(b) {}
    void print() const;
    T1 first;
    T2 second;
};
template <typename T1, typename T2>  // hors classe :
void Pair<T1, T2>::print() const {   //   tout répéter
    std::cout << first << " : " << second << '\n';
}
Pair<int, std::string> p(1, "un");

template <typename T, int N>  // paramètre valeur
class Array { T data[N]; };
template <>                   // spécialisation
class Array<bool, 8> { uint8_t bits; };
// type qui dépend d'un paramètre : typename devant
typename std::vector<T>::iterator it;
\end{code}
\piege{tout le code d'un template va dans le \V|.hpp| (instancié à la compilation),
sinon \V|undefined reference| à l'édition de liens.}
\end{bloc}

\begin{bloc}{Fonctions comme objets}{ctl}
\begin{code}
struct Greater {              // foncteur
    int thr;
    bool operator()(int v) const { return v > thr; }
};
Greater g{42};  g(50);        // true

int thr = 42;
auto f = [thr](int v) { return v > thr; };
auto h = [&](int v) -> bool { return v > thr; };
std::function<bool(int)> fn = f;   // <functional>
\end{code}
\begin{tab3}[0mm]{13mm}
 & \V|[]|          & ne capture rien \\
 & \V|[=]| \V|[&]| & tout par copie · tout par référence \\
 & \V|[x, &y]|     & x copié, y par référence \\
 & \V|[this]|      & accès aux membres de l'objet
\end{tab3}
\note{Lambda = foncteur anonyme (intégrable, contrairement à un pointeur de fonction).
\V|[&]| : danger si elle survit au bloc.}
\end{bloc}

\begin{bloc}{Algorithmes}{lib}
\begin{code}
std::sort(v.begin(), v.end());        // <algorithm>
std::sort(v.begin(), v.end(),
          [](int a, int b) { return a > b; });
auto it = std::find(v.begin(), v.end(), 3);
int n = std::count_if(v.begin(), v.end(), pred);
v.erase(std::remove_if(v.begin(), v.end(), pred),
        v.end());                     // supprime si pred
\end{code}
\end{bloc}
\end{colonne}\colbreak
% ======================= colonne 3 : bibliothèque standard =======================
\begin{colonne}[bibliothèque · move][col-c]
\begin{bloc}{Flux et chaînes}{io}
\begin{code}
std::cout << "x = " << x << '\n';   // sans flush
std::cerr << "erreur" << std::endl; // '\n' + flush
std::cin >> a >> b;                 // mot par mot
std::getline(std::cin, ligne);      // ligne entière
std::istringstream(ligne) >> mot >> n;
std::ifstream in("f.txt");          // ofstream : écrire
while (std::getline(in, ligne)) { ... }
std::cout << std::setw(8) << std::setprecision(3) << x;
\end{code}
\sub{std::string}
\begin{grille}{3}
\V|s.size()|  & \V|s.empty()|   & \V|s[i]  s.at(i)| \\
\V|s + t|     & \V|s += t|      & \V|s == t  s < t| \\
\V|s.find(t)| & \V|s.c_str()|   & \V|s.substr(pos, n)| \\
\V|std::stoi(s)| & \V|std::stod(s)| & \V|std::to_string(x)|
\end{grille}
\note{\V|find| : \V|std::string::npos| si absent · \V|c_str()| : \V|const char *|}
\end{bloc}

\begin{bloc}{Conteneurs}{lib}
\begin{code}
std::vector<int> v = {3, 1, 2};
v.push_back(4);  v.pop_back();  v.size();
v[i];  v.at(i);               // at : out_of_range
v.insert(v.begin() + 1, 9);  v.erase(v.begin());
std::map<std::string, int> m;
m["bob"] = 3;                 // crée la clé si absente
if (m.count("ann")) { ... }   // ou m.find(k) != m.end()
for (auto &[k, val] : m) ...  // C++17, trié par clé
for (auto it = v.begin(); it != v.end(); ++it) *it;
\end{code}
\begin{tab3}[19mm]{16mm}
\V|vector|        & accès O(1)  & contigu, insertion O(n) \\
\V|list|          & accès O(n)  & chaînée, insertion O(1) \\
\V|map| \V|set|   & O(log n)    & arbre trié par clé \\
\V|unordered_map| & O(1) moy.   & table de hachage
\end{tab3}
\piege{\V|push_back| peut réallouer : itérateurs, pointeurs et références vers un
\V|vector| deviennent invalides.}
\end{bloc}

\begin{bloc}{Déplacement (C++11)}{lib}
\begin{schema}[x=1mm,y=1mm,font=\fontcode]
  \node[draw=trait,fill=codebg,inner sep=1pt] (a) at (0,4) {a.data};
  \node[draw=trait,fill=codebg,inner sep=1pt] (b) at (0,0) {b.data};
  \node[draw=cat,line width=0.8pt,inner sep=1.2pt] (t) at (24,2) {tampon};
  \draw[->,dashed,draw=gutter] (a) -- (t);
  \draw[->,line width=0.8pt,draw=cat] (b) -- (t);
  \node[anchor=west,text=note] at (33,2) {b = std::move(a)};
\end{schema}
\begin{code}
Buf(Buf &&o) noexcept       // déplacement : vole
    : data(o.data), n(o.n) { o.data = nullptr; }
Buf b = std::move(a);       // a reste valide, vide
\end{code}
\note{\V|std::move| convertit seulement en \V|T&&|. \textbf{Règle des 5} : dtor, 2 copies,
2 déplacements ; \textbf{des 0} : n'en écrire aucun (\V|vector|, \V|unique_ptr|).}
\end{bloc}

\begin{bloc}{Exceptions}{ctl}
\begin{code}
if (b == 0) throw std::invalid_argument("b nul");
try {
    r = safeDiv(a, b);
} catch (const std::invalid_argument &e) {
    std::cerr << e.what() << '\n';
} catch (const std::exception &e) {   // plus général
} catch (...) { throw; }              // tout : relance
\end{code}
\note{\V|exception| $\leftarrow$ \V|logic_error| (\V|invalid_argument|, \V|out_of_range|),
\V|runtime_error|, \V|bad_alloc|. Locaux détruits à la remontée.}
\end{bloc}
\end{colonne}
\end{document}
